\documentclass[11pt]{article}
\usepackage[margin=1in]{geometry}
\usepackage{amsmath,amssymb,amsthm,mathtools}
\usepackage{enumitem,booktabs,microtype}
\usepackage[colorlinks=true,linkcolor=blue,urlcolor=blue,citecolor=blue]{hyperref}
\setlength{\emergencystretch}{2em}
\newtheorem{theorem}{Theorem}[section]
\newtheorem{proposition}[theorem]{Proposition}
\newtheorem{corollary}[theorem]{Corollary}
\theoremstyle{definition}
\newtheorem{definition}[theorem]{Definition}
\newcommand{\C}{\mathbb C}
\newcommand{\R}{\mathbb R}
\newcommand{\Z}{\mathbb Z}
\newcommand{\HH}{\mathcal H}
\newcommand{\rank}{\operatorname{rank}}
\newcommand{\im}{\operatorname{im}}
\newcommand{\Tr}{\operatorname{Tr}}
\newcommand{\diag}{\operatorname{diag}}
\title{Inverse Simplicial Operator Realization\\[4pt]
\large A companion analytical derivation}
\author{}
\date{}
\hypersetup{pdftitle={Inverse Simplicial Operator Realization}}
\begin{document}
\maketitle
\begin{abstract}
Given a linear operator and declared boundary state spaces, can scalar
simplicial geometry realize that operator for every input? We formulate
the inverse problem in terms of whole harmonic kernels and constrained
boundary energies. We prove an exact boundary normal form, including
singular internal elimination, and derive a composition law. Symbolic
two-channel examples separate exact decorated-geometry representations,
scalar graph phase transport, and attenuated Dirichlet mixing. A
parameterized triangulation program and explicit obstructions identify
the remaining realization problem. No universal scalar triangulation
theorem, matrix-valued edge connection, or physical dynamics is assumed.
\end{abstract}

\section*{Status and relationship to the reading}
This is a companion to
\href{cobordism_functor_reading.pdf}{\emph{The Cobordism Functor, the Theta
Register, and Root Systems}}, whose source is
\path{docs/design/cobordism_functor_reading.tex} \cite{reading}.
That reading remains authoritative for the representation experiment.
This document develops the optional inverse-realization question without
changing its acceptance criteria.

The older \path{bulk_synthesis_theory.tex} discusses entanglement-sourced
sampling. Here we neither reproduce that sampling program nor infer an
operator from the successful extension of one prescribed state.
Whole-operator identities, not sampling convergence, are our criterion.
No engine code or existing document is changed by this companion.

Results marked as theorems or propositions are proved under their stated
hypotheses. Scalar graph examples are exact but are not three-dimensional
manifold fixtures. The final construction program is conditional:
the required simplicial mixing family has not been constructed here.
\clearpage
\tableofcontents
\clearpage

\section{Types, data, and the inverse question}\label{sec:types}
Write $D$ for geometric dimension, $p$ for cochain degree, $g$ for Jacobian
dimension, and $k$ for quantization level. The input and output dimensions
are $n$ and $m$; neither denotes simplex dimension.

\begin{definition}[Admissible geometric data]
A candidate consists of a finite oriented simplicial complex $\Gamma$,
marked input and output subcomplexes, scalar edge lengths $\ell$, and,
where the chosen model permits them, scalar phases $U_e\in U(1)$.
Reversing an edge replaces $U_e$ by $U_e^{-1}=\overline{U_e}$.
It also includes a specified assembly rule, fixed port conventions, and
the allowed boundary conditions. Denote the continuous parameters by
$\lambda$. No freely specified matrix connection or freely specified
cochain mass matrix is an admissible replacement for these data.
\end{definition}

We distinguish three different constructions.
\begin{enumerate}[nosep]
\item \emph{Whole harmonic realization.} States are readouts of the
complete kernel of an assembled harmonic stiffness matrix.
\item \emph{Dirichlet response.} Inputs are externally fixed; remaining
variables minimize energy. Outputs need not extend to a global zero mode.
\item \emph{Theta representation.} States are sections of a quantization
line bundle; operators act through the geometric translation algebra.
Here $\dim\mathcal Q_k=k^g$, not the dimension of $H^1$.
\end{enumerate}
These are not interchangeable. In particular, a result for scalar vertex
potentials ($p=0$) is not a result for harmonic one-cochains or theta
sections. Scalar phases representing a neutral cochain configuration
also do not automatically become a twist of its coboundary operator.

Let $T:\C^n\to\C^m$ be the target and fix positive Hermitian port metrics
$G_{\rm in}$ and $G_{\rm out}$ when making quantum claims. A requested
state pair must satisfy $y=Tx$, or $[y]=[Tx]$ if only rays are prescribed
and $Tx\ne0$. State, operator, and output cannot be prescribed
independently when they violate this equation. The zero operator is
allowed but has no normalized operator-vector state.

\begin{proposition}[Symbolic input determines the operator]\label{prop:coeff}
For linear $F,T:\C^n\to\C^m$, equality
$F(\sum_j x_j e_j)=T(\sum_j x_j e_j)$ for independent complex coefficients
$x_j$ holds if and only if $F=T$.
\end{proposition}
\begin{proof}
Set $x_j=1$ and all other coefficients to zero to obtain equality of
every column. Conversely, equality of columns implies the identity by
linearity.
\end{proof}
Fitting a single nonzero input only constrains one column in an adapted
basis: every $T+N$ with $Nx=0$ gives the same output. Fitting a normalized
ray additionally loses amplitude. No target-dependent rebasing or output
renormalization is permitted to turn such a fit into an operator proof.

\section{What is already invertible analytically}\label{sec:symbol}
The reading's reduced geometric translation algebra has basis
$e_{a,b}$ for $a,b\in(\Z/k)^g$, with
\[
 e_{a,b}e_{a',b'}=\omega^{b\cdot a'}e_{a+a',b+b'},\qquad
 \omega=e^{2\pi i/k}.
\]
Its quantization sends $e_{a,b}$ to $D_{a,b}=X^aZ^b$ on a
$d=k^g$ dimensional theta register. Orthogonality gives the exact inverse
\begin{equation}\label{eq:symbol}
 s_T=\frac1d\sum_{a,b}\Tr(D_{a,b}^\dagger T)e_{a,b},
 \qquad \mathfrak Q(s_T)=T.
\end{equation}
This is an inverse into \emph{complex linear combinations of decorated
geometric operations}, not into one undecorated triangulation.
Its curve interpretation is developed by Gelca--Uribe \cite{gelca}.
Coefficients in \eqref{eq:symbol} retain the information in $T$; they
are not proved to arise from edge lengths.

At $k=2,g=1$, put
\[
 X=\begin{pmatrix}0&1\\1&0\end{pmatrix},\qquad
 Z=\begin{pmatrix}1&0\\0&-1\end{pmatrix}.
\]
For $T=\begin{psmallmatrix}a&b\\c&d\end{psmallmatrix}$ the inverse is
\begin{equation}\label{eq:qubit-symbol}
 s_T=\tfrac{a+d}{2}e_{0,0}
     +\tfrac{b+c}{2}e_{1,0}
     +\tfrac{a-d}{2}e_{0,1}
     +\tfrac{c-b}{2}e_{1,1}.
\end{equation}
Indeed $XZ=\begin{psmallmatrix}0&-1\\1&0\end{psmallmatrix}$, so all four
entries are recovered. For example
\[
 U(\theta)=
 \begin{pmatrix}\cos\theta&-\sin\theta\\
                 \sin\theta&\cos\theta\end{pmatrix}
 \quad\longleftrightarrow\quad
 \cos\theta\,e_{0,0}+\sin\theta\,e_{1,1}.
\]
This is an exact continuous mixing family in the \emph{linearized
geometric algebra}. It is not an assertion that a cylinder's native
period transport varies continuously with its lengths.
For rectangular operators use matrix units between distinct registers;
do not use a square Weyl basis of the wrong dimension.

\section{Assembly from scalar simplices}\label{sec:assembly}
For the positive-real metric model, let $M_p(\ell)>0$ be the assembled
Whitney cochain mass and let $d_p$ be the coboundary. Define
\[
 \delta_p=M_{p-1}^{-1}d_{p-1}^\dagger M_p,\qquad
 L_p=d_{p-1}\delta_p+\delta_{p+1}d_p.
\]
The Hermitian stiffness is
\begin{equation}\label{eq:stiffness}
 K_p=M_pL_p
 =M_pd_{p-1}M_{p-1}^{-1}d_{p-1}^\dagger M_p
   +d_p^\dagger M_{p+1}d_p.
\end{equation}
At an endpoint degree the absent term is zero. Its energy is
\[
 h^\dagger K_ph
 =\|\delta_ph\|_{M_{p-1}}^2+\|d_ph\|_{M_{p+1}}^2,
\]
so $K_p\ge0$ and $\ker K_p=\ker L_p$. A flat unitary local system
admits the analogous construction only with consistently assembled
twisted maps and a verified $d_{p+1}^Ud_p^U=0$.
Arbitrary phases on filled triangles need not satisfy this condition.

For a simplex with vertices $v_0,\ldots,v_D$, its local Euclidean Gram is
\begin{equation}\label{eq:localgram}
 Q_{ab}=\frac{\ell_{0a}^2+\ell_{0b}^2-\ell_{ab}^2}{2},
 \qquad 1\le a,b\le D.
\end{equation}
The condition $Q>0$ gives a nondegenerate Euclidean simplex.
Lengths on shared faces must agree. This is an intrinsic piecewise-flat
metric; it does not assert a global isometric embedding in $\R^D$.
The masses in \eqref{eq:stiffness} must be assembled from these local
metrics, not chosen afterward to force $T$.

Incidence matrices satisfy $\partial_p\partial_{p+1}=0$.
That identity alone does not characterize a simplicial manifold:
simplex membership, face closure, orientation, boundary identifications,
and the required sphere or ball links must also hold.
A tetrahedron requires all four compatible triangular faces and its
three-dimensional interior. Filling every possible clique can destroy
the intended cohomology.

\paragraph{Signature boundary.}
Complex-bilinear or Lorentzian pencils need separate kernel and rank
certificates. Neither the positive energy arguments below nor a
positive quantum norm follows from those pencils automatically.
The algebraic whole-kernel criterion in the next section still applies
to a computed geometric kernel, without assuming positivity.

\section{The exact whole-kernel inverse equations}\label{sec:kernel}
Let $Z\in\C^{N\times r}$ be a full basis of $\ker K_\Gamma(\lambda)$.
The declared linear readouts are $R_{\rm in}$ and $R_{\rm out}$:
\[
 A=R_{\rm in}Z\in\C^{n\times r},\qquad
 B=R_{\rm out}Z\in\C^{m\times r}.
\]
Readout rules may depend on geometry in a prescribed way but not on the
target chosen during recovery. Period readouts and whole-metric Gram
readouts are distinct. Use the reading's native dual/primal conventions
and existing \texttt{PencilSchur.gramBlock}, not a new mass or an
unjustified Hermitian replacement for a bilinear pairing.

\begin{theorem}[Realization on all inputs]\label{thm:kernel}
The complete boundary relation is the graph of a specified operator $T$
if and only if
\begin{equation}\label{eq:kernel}
 \rank A=n,\qquad B=TA.
\end{equation}
Equivalently, the graph exists for some unique $T$ precisely when
$A$ is onto and $\ker A\subseteq\ker B$.
\end{theorem}
\begin{proof}
Every bulk zero mode has the form $Zc$. If \eqref{eq:kernel} holds,
every $x$ equals $Ac$ for some $c$, and its output is $Bc=Tx$.
If two extensions differ, their difference lies in $\ker A$ and is
annihilated by $B$. Conversely, existence of every input implies that
$A$ is onto, and the specified graph forces $Bc=TAc$ for all $c$.
For the equivalent criterion choose a right inverse $C$ of $A$.
Since $c-CAc\in\ker A$, one has $B=BCA$ and $T=BC$.
Any other right inverse differs by a map into $\ker A$, so gives
the same $T$.
\end{proof}

\begin{corollary}[A symbolic formulation without a chosen kernel basis]
It suffices to find $V\in\C^{N\times n}$ satisfying
\[
 K_\Gamma V=0,\qquad R_{\rm in}V=I,\qquad R_{\rm out}V=T,
\]
\emph{and} certify
$\ker K_\Gamma\cap\ker R_{\rm in}\subseteq\ker R_{\rm out}$.
\end{corollary}
The last condition cannot be omitted: $V$ alone supplies witnesses but
does not rule out additional visible output modes. Invisible bulk modes
are allowed. Full-rank minors and a complete null-space argument are
analytical certificates; numerical ranks require explicit tolerances,
gaps, and conditioning.

If a native engine frame is a chain frame $\Phi$ while periods act on
cochain images $Z=G^U\Phi$, perform that conversion before applying
\eqref{eq:kernel}. Equal dimensions do not certify equal spaces.

\section{Constrained boundary energy, including null modes}\label{sec:schur}
To use an effective boundary energy, collect the declared port
coordinates in a map
\[
 P=\begin{pmatrix}R_{\rm in}\\R_{\rm out}\end{pmatrix}:
 \C^N\longrightarrow\C^{n+m}.
\]
Assume $P$ is onto on the ambient cochain space, not necessarily on the
harmonic space. Choose a right inverse $E$ of $P$ and a full column basis
$F$ of $\ker P$. Every field with port values $b=(x,y)$ has the unique form
$h=Eb+Fz$. Thus
\[
 h^\dagger Kh=
 \begin{pmatrix}b\\z\end{pmatrix}^{\!\dagger}
 \begin{pmatrix}K_{bb}&K_{bz}\\K_{zb}&K_{zz}\end{pmatrix}
 \begin{pmatrix}b\\z\end{pmatrix},
\]
where the blocks are obtained by congruence with $[E,F]$.
This accounts for all unobserved directions, including unretained
boundary modes. A principal submatrix in arbitrarily chosen boundary
coordinates need not be the effective matrix for the declared ports.
If $P$ is not onto, retain its feasibility constraints explicitly.

\begin{theorem}[Positive singular elimination]\label{thm:schur}
For a positive semidefinite block matrix as above,
$\im K_{zb}\subseteq\im K_{zz}$. The constrained minimum exists for all $b$
and equals
\begin{equation}\label{eq:schur}
 \mathcal E(b)=\min_{Ph=b}h^\dagger Kh=b^\dagger S_\Gamma b,
 \qquad
 S_\Gamma=K_{bb}-K_{bz}K_{zz}^{+}K_{zb}\ge0.
\end{equation}
Here $+$ is the Moore--Penrose inverse in these declared coordinates.
Moreover $\ker S_\Gamma=P(\ker K)$.
\end{theorem}
\begin{proof}
If $v\in\ker K_{zz}$, positivity implies that the block quadratic
form at $(b,tv)$ cannot have a nonzero term linear in arbitrary complex
$t$; therefore $K_{bz}v=0$. This gives the range inclusion.
The minimizing internal variables are
$z=-K_{zz}^{+}K_{zb}b+v$ with $v\in\ker K_{zz}$.
Substitution gives \eqref{eq:schur}, which is nonnegative.
If $b\in\ker S_\Gamma$, the reconstructed minimizing field has zero
energy and hence lies in $\ker K$, since $K\ge0$.
The converse follows by inserting any zero-energy field with ports $b$.
\end{proof}
This use of a pseudoinverse is justified by positivity and the range
proof; it is not an unchecked replacement for a singular inverse.
The minimized quadratic form is independent of the choice of $E,F$.
For a general complex non-Hermitian pencil, compatibility must instead
be established directly. Schur reduction is elimination, not a partial
trace. Graph-Laplacian and simplex versions of this relationship have
independent precedents \cite{devriendt}.

\section{A complete boundary normal form}\label{sec:normal}
\begin{theorem}[Graph-kernel normal form]\label{thm:normal}
A Hermitian positive semidefinite matrix $S$ on $\C^n\oplus\C^m$
has kernel exactly $\{(x,Tx):x\in\C^n\}$ if and only if
\begin{equation}\label{eq:normal}
 S=[-T,I]^\dagger W[-T,I]
 =\begin{pmatrix}T^\dagger WT&-T^\dagger W\\-WT&W\end{pmatrix}
 \quad\text{for some }W>0.
\end{equation}
\end{theorem}
\begin{proof}
For the forward direction write $W=S_{yy}$. For nonzero $y$, $(0,y)$
does not lie in the stated kernel, so $y^\dagger Wy>0$.
The identity $S\binom{I}{T}=0$ gives
$S_{yx}=-WT$ and $S_{xx}=T^\dagger WT$; Hermitian symmetry gives
$S_{xy}=-T^\dagger W$. Conversely, \eqref{eq:normal} has energy
$(y-Tx)^\dagger W(y-Tx)$, which vanishes exactly on that graph.
\end{proof}

For the positive whole-kernel model with onto ambient port map, the
inverse problem is therefore exactly
\begin{equation}\label{eq:inverse}
 S_\Gamma(\lambda)=[-T,I]^\dagger W[-T,I],
 \qquad W>0,\quad(\Gamma,\lambda)\text{ admissible}.
\end{equation}
Leaving $W$ unknown does not weaken the operator condition: it permits
different costs for deviations from the same graph.
For $n=m=2$ and $W=I$, a symbolic
$T=\begin{psmallmatrix}a&b\\c&d\end{psmallmatrix}$ requires
\[
 S_{xx}=
 \begin{pmatrix}|a|^2+|c|^2&\bar a b+\bar c d\\
                \bar b a+\bar d c&|b|^2+|d|^2\end{pmatrix},
 \quad S_{yx}=-T,\quad S_{yy}=I.
\]
Thus the backward calculation specifies blocks of an effective
quadratic form; scalar simplex assembly must realize those blocks
simultaneously. Choosing $S$ directly from $T$ verifies the algebra
but does not solve \eqref{eq:inverse}.

\section{Two exact scalar graph controls}\label{sec:controls}
These examples state precisely what elementary connectivity can do.
They are one-dimensional simplicial complexes with scalar phases and
fixed vertex ports, not claims about a three-dimensional Hodge register.

\subsection{One-to-one phase transport}
Take $n$ disjoint edges joining $x_j$ to $y_j$, with $w_j>0$ and
scalar phase $u_j=e^{i\phi_j}$. Their connection energy is
\[
 \mathcal E(x,y)=\sum_{j=1}^n w_j|y_j-u_jx_j|^2.
\]
It has the normal form \eqref{eq:normal} with
$T=\diag(u_1,\ldots,u_n)$ and $W=\diag(w_1,\ldots,w_n)$.
For $\phi_1=-\theta/2,\phi_2=\theta/2$ it realizes the two-channel
phase operator $\diag(e^{-i\theta/2},e^{i\theta/2})$ exactly.
Each component has one independent zero mode.

For the ordinary scalar energy on a Euclidean interval,
$w_j=1/\ell_j$ with unit material coefficient, so its positive weight
comes from a length. The prescribed $u_j$ is still separate scalar
phase data. On a tree that phase is pure gauge; the displayed operator
has meaning only relative to fixed input/output trivializations.
This is not a gauge-independent phase generated by the length.
Connecting all channels into a single positive degree-zero graph
changes the zero-mode count, as shown below.

\subsection{Mixing by Dirichlet response, with its exact attenuation}
Connect each of two input vertices to each of two output vertices,
with unit weights and three phases $+1$ and one phase $-1$:
\begin{equation}\label{eq:mixenergy}
 \mathcal E=
 |y_1-x_1|^2+|y_1-x_2|^2+
 |y_2-x_1|^2+|y_2+x_2|^2.
\end{equation}
Hold $x$ fixed and vary $y$. Stationarity gives
\[
 2y_1=x_1+x_2,\qquad 2y_2=x_1-x_2,\qquad
 y=\frac{H}{\sqrt2}x,\quad
 H=\frac1{\sqrt2}\begin{pmatrix}1&1\\1&-1\end{pmatrix}.
\]
This is an exact scalar construction of the response $H/\sqrt2$,
not of $H$ with its amplitude preserved. Indeed
$\|y\|^2=\|x\|^2/2$, and the minimized total energy is $\|x\|^2$.
The whole graph has no nonzero zero mode: zero energy would require
simultaneously $y_1=x_1=x_2$ and $y_2=x_1=-x_2$.

Its phase holonomy around the four-cycle is $-1$. Filling that cycle
with a flat triangular disk is impossible without changing phase data
or introducing curvature. Merely adding tetrahedra is not a
phase-preserving upgrade of this example.
As a quantum Kraus operator, $H/\sqrt2$ has success probability $1/2$
on a normalized input; a completion is needed for a channel.
Renormalizing its output reproduces the Hadamard ray but loses that
probability. A classical Dirichlet graph alone is not a quantum
implementation of that measurement.

\begin{proposition}[General attenuated scalar response]\label{prop:attenuation}
Let $T\in\C^{m\times n}$ and choose
$c>0$ with $c\ge\max_j\sum_i|T_{ji}|$.
There is a positive scalar phase graph, with fixed input and grounded
vertices, whose output Dirichlet response is $y=Tx/c$.
\end{proposition}
\begin{proof}
For nonzero $T_{ji}$ join $x_i$ to $y_j$ with weight $|T_{ji}|$
and phase $T_{ji}/|T_{ji}|$. Omit zero entries. Connect $y_j$ to a fixed
zero vertex with weight $c-\sum_i|T_{ji}|$ when this is positive.
Each output stationarity equation is $cy_j=\sum_iT_{ji}x_i$.
Positive weights can be encoded by interval lengths $1/w$.
\end{proof}
This proves a scalar graph response statement, not a whole-kernel,
manifold, or theta synthesis theorem. Downstream loading can change
a Dirichlet response, so this control must not be composed by assuming
that each output is an independently buffered input.
Multiplication by $c$ at readout is an explicit gain, not a geometric
unitarity proof with fixed metrics.

\section{Obstructions and what increasing dimension does not prove}
\label{sec:obstructions}
\begin{proposition}[Positive scalar zero-mode bound]\label{prop:zerobound}
For a connected graph with energy
$\sum_{(i,j)}w_{ij}|f_j-u_{ij}f_i|^2$, $w_{ij}>0$ and $|u_{ij}|=1$,
the zero-mode dimension is at most one. It is one precisely when all
cycle holonomies are trivial.
\end{proposition}
\begin{proof}
Zero energy forces $f_j=u_{ij}f_i$ on every edge. A spanning tree
determines every value from one root value. Each remaining edge either
agrees or forces the root value to zero. Agreement is exactly trivial
cycle holonomy.
\end{proof}
More vertices in a connected degree-zero graph do not create an
arbitrary-dimensional zero-energy input register. Higher cochain degree
can avoid this particular bound, but requires a separate construction.

\begin{proposition}[Dirichlet maximum bound]\label{prop:maxbound}
In a positive scalar phase graph with a nonempty fixed boundary meeting
every connected component, the free-vertex Dirichlet solution satisfies
$|f_v|\le\max_{\partial}|f|$.
Consequently its boundary-to-output matrix $F$ obeys
$\sum_i|F_{ji}|\le1$ for each output row.
\end{proposition}
\begin{proof}
Stationarity expresses each free value as a convex weighted average of
phase-rotated neighboring values. An interior modulus maximum larger
than all boundary moduli would propagate along positive edges to the
boundary, a contradiction. For a fixed output row choose independent
boundary phases so that all nonzero terms $F_{ji}x_i$ align, with
$|x_i|=1$. The modulus bound then gives the row inequality.
\end{proof}
This bound excludes the unattenuated Hadamard operator at raw vertex
ports, whose row absolute sums are $\sqrt2$. It does not exclude
Hadamard in a theta register, a modal readout, or a different dynamical
model.

\paragraph{Topological period obstruction.}
For ordinary untwisted cohomology with fixed markings, the boundary
relation is independent of lengths. Over rational period coordinates
it is defined over $\mathbb Q$; where it is a graph its matrix is rational.
Integral lattice isomorphisms give unimodular integral maps. For an
oriented manifold, the applicable boundary intersection pairing further
constrains the relation. Hence varying lengths alone cannot generate a
continuous family of arbitrary complex period operators.
Changing the local system is not a metric deformation with that local
system held fixed.

\paragraph{Quantum versus topological dimension.}
A product-polarized theta register has dimension $k^g$ on a rank-$2g$
period lattice. Replacing this with a $k^g$ dimensional harmonic kernel
is a new model, not a change of notation. Similarly, a simplicial join
shifts homology degree; it does not automatically implement the reading's
theta tensor product or its dual-input pairing.

\paragraph{Parameter count and nonuniqueness.}
On a smooth constant-rank parameter region, a family with $r$ effective
real parameters cannot cover an open subset of a $q$ dimensional target
manifold if $r<q$. General complex $m$ by $n$ matrices have $q=2mn$;
unitary $n$ by $n$ matrices have $q=n^2$ (one fewer projectively).
This is a necessary dimension bound, not a sufficient realization test.
A deficient derivative at one point alone is not a proof of failure:
$x\mapsto x^3$ is locally onto at zero despite zero derivative.
Different interiors can have identical responses. Even scalar
conductances $w_1,w_2$ in series reduce to $w_1w_2/(w_1+w_2)$, so an
inverse should return a family or a stated normal form, not claim
unique connectivity.

\section{Gluing: the compositional theorem and its hypotheses}\label{sec:gluing}
\begin{theorem}[Composition of positive graph constraints]\label{thm:gluing}
Let two effective boundary energies be
\[
 E_1=(y-T_1x)^\dagger W_1(y-T_1x),\qquad
 E_2=(z-T_2y)^\dagger W_2(z-T_2y),\quad W_1,W_2>0.
\]
Their sum, minimized over the shared $y$, has the form
\[
 \min_y(E_1+E_2)
 =(z-T_2T_1x)^\dagger W_{\rm eff}(z-T_2T_1x),
\]
where
\begin{equation}\label{eq:weff}
\begin{split}
 W_{\rm eff}
 &=W_2-W_2T_2(W_1+T_2^\dagger W_2T_2)^{-1}T_2^\dagger W_2\\
 &=(W_2^{-1}+T_2W_1^{-1}T_2^\dagger)^{-1}>0.
\end{split}
\end{equation}
\end{theorem}
\begin{proof}
Set $u=y-T_1x$ and $r=z-T_2T_1x$. The energy becomes
$u^\dagger W_1u+(r-T_2u)^\dagger W_2(r-T_2u)$.
Its unique minimizer solves
$(W_1+T_2^\dagger W_2T_2)u=T_2^\dagger W_2r$.
Substitution proves the first expression; multiplication, or the
Woodbury identity, gives the second. Its inverse is positive definite.
\end{proof}

This is an algebraic composition theorem. For actual simplicial gluing
one must prove that the assembled glued problem has this summed energy
or the same boundary relation. Match interface cells, orientations,
phase trivializations and port metrics; retain all interface variables.
Eliminating a variable separately on both sides and identifying only a
few summaries need not give the same result.
In particular the inverse mass in \eqref{eq:stiffness} makes naive
additivity of separately assembled Hodge stiffness matrices unjustified.
The reading's cohomological gluing theorem has its own topological
hypotheses and does not remove this metric-assembly obligation.

A collar between matched surfaces can be triangulated using ordered
prism subdivisions: an edge prism needs a diagonal and a triangle prism
three tetrahedra. Globally consistent orderings ensure matching faces
\cite{hatcher}. This constructs a cylinder topology, not an arbitrary
operator action. Whole three-dimensional complexes crossed with an
interval instead require four-dimensional simplices.

\section{A constructive analytical program}\label{sec:program}
The proposed general form is a finite layered collar with fixed marked
boundary ports. Layers contain a controlled list of internal simplices
and possible couplings; lengths and permitted phases are symbolic.
The list is not assumed universal.

\begin{enumerate}[leftmargin=*,label=\textbf{Step \arabic*.}]
\item \textbf{Fix the semantics.}
Choose harmonic degree, positive or complex formulation, state space,
readouts, gauge conventions and boundary metrics. Freeze the original
input/output basis. State whether the target is an exact linear map,
a projective unitary, a contraction, or a channel branch.

\item \textbf{Choose an admissible connectivity family.}
Build oriented incidence from actual simplices with compatible
interfaces. One-to-one edges alone do not close mixed triangles;
introduce and record the needed diagonals and higher faces.
If manifolds are required, impose the local link conditions.
Keep a finite size bound when claiming exhaustive failure.

\item \textbf{Derive the action on a symbolic input.}
Use the full-kernel equations \eqref{eq:kernel}, or, under the positive
and ambient-port hypotheses, the exact inverse system \eqref{eq:inverse}.
For fixed connectivity solve continuous parameters before increasing
combinatorial complexity. Preserve singular compatibility conditions
when clearing denominators.

\item \textbf{Enforce geometric feasibility.}
Require \eqref{eq:localgram}, shared-face agreement, boundary pinning,
and local-system constraints. A solution for an arbitrary SPD matrix is
not a solution for a Whitney mass. Parameter zeros that remove edges or
collapse simplices lie on different combinatorial or inadmissible strata.

\item \textbf{Solve elementary families and glue.}
The phase and attenuated response examples above are analytical
controls. The missing native ingredient is an exact, continuously
parameterized two-channel mixing family in the \emph{same} state space
and admissible simplicial model, together with an actual assembly/gluing
proof. A scalar vertex response control is not such an ingredient.

\item \textbf{Recover independently and classify the result.}
From the resulting geometry alone, recover the full operator.
Check the entire kernel, all input columns and declared norms.
Report exact representation, exact graph response, exact native
realization, approximation, or an obstruction within the stated family.
A failed numerical search is not an impossibility proof.
\end{enumerate}

For an assembly and readout that are algebraic in squared lengths,
positive volume variables and phase components, the fixed-connectivity
equations can be expressed using real polynomial equalities and
inequalities, auxiliary inverse variables, and rank minors.
Phase components obey $(\Re U_e)^2+(\Im U_e)^2=1$.
This gives an exact feasibility formulation; it does not assert
a closed-form solution or computational tractability.
Theta series or other nonalgebraic readouts require a different
analytic treatment.

\begin{proposition}[Conditional finite-unitary construction]\label{prop:conditional}
Suppose a single declared geometric model realizes every two-channel
unitary on any chosen pair of channels and identity on the spectators,
realizes single-channel phases, and its actual gluing implements
operator composition without leakage or unrecorded scaling.
Then every finite-dimensional unitary on those channels has a finite
geometric realization.
\end{proposition}
\begin{proof}
Successive two-coordinate unitary row operations annihilate entries
below the diagonal of a target unitary. The resulting upper triangular
unitary is diagonal with unit-modulus entries: orthonormality of its
columns proves this inductively. Reversing the row operations and
including the diagonal phases expresses the target as a finite product
of the assumed realizable factors. The gluing hypothesis realizes that
product.
\end{proof}
This is a reduction of the missing theorem, not its proof.
The unitary premise does not cover arbitrary nonunitary matrices.
A contraction can instead be considered with an explicitly specified
ancilla and projection; their geometry, readouts and success
probabilities become additional obligations.

\section{States, norms, and the acceptance certificate}\label{sec:accept}
Once $T$ is realized, all compatible state pairs follow by linearity.
For fixed positive metrics, isometry means
\begin{equation}\label{eq:norm}
 T^\dagger G_{\rm out}T=G_{\rm in}.
\end{equation}
Do not change $G_{\rm out}$ after seeing $T$ to manufacture this identity.
A valid nonunitary linear realization need not satisfy it.
For a physical channel branch require
$T^\dagger G_{\rm out}T\le G_{\rm in}$ after the appropriate
Hilbert-space identification, and specify the other branches for a
trace-preserving channel.

Output-major vectorization belongs to
$\HH_{\rm out}\otimes\overline{\HH_{\rm in}}$.
In orthonormal coordinates its squared norm is $\Tr(T^\dagger T)$.
Keeping only its normalized ray loses scale. With nonorthonormal
coordinates, first use the declared Gram metrics to identify
orthonormal Hilbert coordinates. An operator vector is not an
unexplained cochain on a boundary mesh.

A completed native construction should supply:
\begin{enumerate}[nosep]
\item The complete simplex list or a parameterized construction rule;
lengths, phases, markings, and boundary conditions.
\item A derivation that its assembled matrices are exactly the matrices
used in the proof, with no target-dependent readout.
\item Complete harmonic-kernel and port-rank certificates, including
invisible modes and singular compatibility.
\item A symbolic identity on every input, with exact parameter domain,
not only a state residual or a favorable numerical sample.
\item A separate norm or channel certificate when such behavior is claimed.
\item An actual interface assembly proof before claiming composability.
\end{enumerate}
Numerical substitution checks are useful audits of these identities,
not substitutes for their analytical proofs.

\section{What is established and what remains}\label{sec:verdict}
The exact inverse to a reduced decorated-geometry symbol is already
available. Theorems \ref{thm:kernel}--\ref{thm:normal} identify the precise
scalar-geometric inverse problem, and Theorem~\ref{thm:gluing} states
the energy-level composition law. The scalar phase and attenuated mixing
examples are actual analytical constructions in their stated graph
models. The bounds in Section~\ref{sec:obstructions} explain why they
cannot simply be relabelled as arbitrary harmonic quantum gates.

The outstanding task is an admissible simplicial family solving
\eqref{eq:inverse}, or \eqref{eq:kernel} in the nonpositive setting,
for the desired mixing operators in the chosen register, with
compatible physical norms and actual gluing when claimed.
Higher dimension, richer connectivity, and scalar phases are available
design variables; none alone proves that this family exists.

Thus the answer to the inverse question is constructive but conditional:
we can solve backward to exact boundary constraints and several scalar
controls now. A general triangulation formula requires the missing
geometric realization theorem, not another identification of matrices
with geometric symbols.

\clearpage
\begin{thebibliography}{9}
\bibitem{reading}
\emph{The Cobordism Functor, the Theta Register, and Root Systems}.
Tessera, \path{docs/design/cobordism_functor_reading.tex}.
In particular: \emph{The full geometric operator algebra},
\emph{Frozen whole-kernel recovery}, \emph{Pairing and Gram conventions},
and \emph{Internal elimination and actual gluing}.

\bibitem{gelca}
R. Gelca and A. Uribe,
\emph{From classical theta functions to topological quantum field theory}.
\href{https://arxiv.org/abs/1006.3252}{arXiv:1006.3252}.
Supports the curve-algebra representation, not unrestricted native
simplicial synthesis.

\bibitem{devriendt}
K. Devriendt,
\emph{Effective resistance is more than distance: Laplacians, Simplices
and the Schur complement}.
\href{https://arxiv.org/abs/2010.04521}{arXiv:2010.04521}.
Provides graph/simplex and Schur-reduction context, not the manifold
realization assertion left open here.

\bibitem{hatcher}
A. Hatcher, \emph{Algebraic Topology}, proof of Theorem 2.10.
\href{https://pi.math.cornell.edu/~hatcher/AT/AT.pdf}{Author's edition}.
The ordered triangulation of a simplex times an interval.
\end{thebibliography}
\end{document}
